\documentclass{beamer}
\usepackage[utf8]{inputenc}
\usetheme{Madrid}
\usecolortheme{default}
\usepackage{amsmath,amssymb,amsfonts,amsthm}
\usepackage{txfonts}
\usepackage{tkz-euclide}
\usepackage{listings}
\usepackage{adjustbox}
\usepackage{array}
\usepackage{tabularx}
\usepackage{gvv}
\usepackage{lmodern}
\usepackage{circuitikz}
\usepackage{tikz}
\usepackage{graphicx}
\setbeamertemplate{page number in head/foot}[totalframenumber]
\usepackage{tcolorbox}
\tcbuselibrary{minted,breakable,xparse,skins}

\definecolor{bg}{gray}{0.95}
\DeclareTCBListing{mintedbox}{O{}m!O{}}{%
  breakable=true,
  listing engine=minted,
  listing only,
  minted language=#2,
  minted style=default,
  minted options={%
    linenos,
    gobble=0,
    breaklines=true,
    fontsize=\small,
    numbersep=8pt,
    #1},
  boxsep=0pt,
  left=25pt,
  right=0pt,
  top=3pt,
  bottom=3pt,
  arc=5pt,
  leftrule=0pt,
  rightrule=0pt,
  bottomrule=2pt,
  toprule=2pt,
  colback=bg,
  colframe=orange!70,
  enhanced,
  overlay={%
    \begin{tcbclipinterior}
    \fill[orange!20!white] (frame.south west) rectangle ([xshift=20pt]frame.north west);
    \end{tcbclipinterior}},
  #3,
}
\lstset{
    language=C,
    basicstyle=\ttfamily\small,
    keywordstyle=\color{blue},
    stringstyle=\color{orange},
    commentstyle=\color{green!60!black},
    numbers=left,
    numberstyle=\tiny\color{gray},
    breaklines=true,
    showstringspaces=false,
}

\title{12.632}
\author{AI25BTECH11034 - Sujal Chauhan}

\begin{document}

\frame{\titlepage}

\begin{frame}{Problem}
12.632\\
For given $3\times3$ matrix $\Vec{A}$ eigonvalues are 1,2,3 then value of $trace(\Vec{A}^{-1})$ is.

\end{frame}
\begin{frame}{Theory}


\section*{Theorem: Eigenvalues and Eigenvectors of the Inverse Matrix}

Let \(A\) be an invertible \(n \times n\) matrix with eigenvalue \(\lambda\) and corresponding eigenvector \(\mathbf{v}\), i.e.,
\[
A \mathbf{v} = \lambda \mathbf{v}
\]
where \(\lambda \neq 0\) and \(\mathbf{v} \neq \mathbf{0}.\)

\subsection*{Claim}

\begin{itemize}
    \item The eigenvectors of \(A^{-1}\) are the same as those of \(A\).
    \item The eigenvalues of \(A^{-1}\) are the reciprocals of the eigenvalues of \(A\), i.e., \(\frac{1}{\lambda}\).
\end{itemize}

\subsection*{Proof}

Starting from the eigenvalue equation:
\[
A \mathbf{v} = \lambda \mathbf{v}
\]
\end{frame}
\begin{frame}{Theory}
    
Since \(A\) is invertible and \(\lambda \neq 0\), multiply both sides by \(A^{-1}\):
\[
A^{-1} A \mathbf{v} = A^{-1} \lambda \mathbf{v}
\]
\[
\mathbf{v} = \lambda A^{-1} \mathbf{v}
\]

Rearranging:
\[
A^{-1} \mathbf{v} = \frac{1}{\lambda} \mathbf{v}
\]

This shows that \(\mathbf{v}\) is an eigenvector of \(A^{-1}\) with eigenvalue \(\frac{1}{\lambda}\).

\subsection*{Conclusion}

The eigenvalues of \(A^{-1}\) are the reciprocals of the eigenvalues of \(A\), and both matrices share the same eigenvectors.
\end{frame}

\begin{frame}{Theory}
    

\section*{Proof: Sum of Eigenvalues Equals Trace Using Diagonalization}

Let \(A\) be an \(n \times n\) diagonalizable matrix. Then there exists an invertible matrix \(P\) and a diagonal matrix \(D\) such that
\[
A = P D P^{-1}
\]
where
\[
D = \mathrm{diag}(\lambda_1, \lambda_2, \ldots, \lambda_n)
\]
and \(\lambda_1, \lambda_2, \ldots, \lambda_n\) are the eigenvalues of \(A\).

\bigskip

Taking the trace on both sides, we have
\[
\mathrm{tr}(A) = \mathrm{tr}(P D P^{-1}).
\]
\end{frame}
\begin{frame}[Theory]
Using the cyclic property of trace, which states that \(\mathrm{tr}(XYZ) = \mathrm{tr}(ZXY) = \mathrm{tr}(YZX)\) for any matrices \(X, Y, Z\) where the product is defined, we get
\[
\mathrm{tr}(P D P^{-1}) = \mathrm{tr}(D P^{-1} P) = \mathrm{tr}(D I) = \mathrm{tr}(D).
\]

Since \(D\) is diagonal, its trace is the sum of its diagonal entries, i.e., the sum of eigenvalues of \(A\):
\[
\mathrm{tr}(D) = \lambda_1 + \lambda_2 + \cdots + \lambda_n.
\]

Therefore,
\[
\boxed{
\mathrm{tr}(A) = \sum_{i=1}^n \lambda_i.
}
\]

This completes the proof.



\end{frame}
\begin{frame}{Theory}
  
\section*{Proof: Cyclic Property of Trace}

Let \(A, B, C\) be \(n \times n\) matrices. The trace of the product \(ABC\) is defined as the sum of the diagonal elements:
\[
\mathrm{tr}(ABC) = \sum_{i=1}^n (ABC)_{ii}.
\]

\bigskip

Recall that matrix multiplication components are given by
\[
(ABC)_{ii} = \sum_{j=1}^n (AB)_{ij} C_{ji} = \sum_{j=1}^n \sum_{k=1}^n A_{ik} B_{kj} C_{ji}.
\]
\end{frame}
\begin{frame}{Theory}
Thus,
\[
\mathrm{tr}(ABC) = \sum_{i=1}^n \sum_{j=1}^n \sum_{k=1}^n A_{ik} B_{kj} C_{ji}.
\]

\bigskip

By changing the order of summation and renaming indices, we can write:
\[
\mathrm{tr}(BCA) = \sum_{j=1}^n \sum_{k=1}^n \sum_{i=1}^n B_{kj} C_{ji} A_{ik} = \sum_{i=1}^n \sum_{j=1}^n \sum_{k=1}^n A_{ik} B_{kj} C_{ji} = \mathrm{tr}(ABC).
\]

Similarly,
\[
\mathrm{tr}(CAB) = \sum_{k=1}^n \sum_{i=1}^n \sum_{j=1}^n C_{ji} A_{ik} B_{kj} = \sum_{i=1}^n \sum_{j=1}^n \sum_{k=1}^n A_{ik} B_{kj} C_{ji} = \mathrm{tr}(ABC).
\]

\bigskip
\end{frame}
\begin{frame}{Theory}
Hence, we have shown that
\[
\mathrm{tr}(ABC) = \mathrm{tr}(BCA) = \mathrm{tr}(CAB).
\]

This proves the cyclic property of trace for three matrices. The property extends similarly to any product of matrices where multiplication is defined.



\end{frame}

\begin{frame}{Solution}
    Now using properties mentioned above :
    eigon values of $\Vec{A}^{-1}$ are $\frac{1}{1},\frac{1}{2},\frac{1}{3}$
    
\end{frame}    
\begin{frame}{Solution}
   \begin{align}
       trace(\Vec{A}^{-1})= \frac{1}{1}+\frac{1}{2}+\frac{1}{3}=\frac{11}{6}
   \end{align} 
\end{frame}

\end{document}
